% GENERATED FILE. Edit canonical lessons, then run npm run generate:books.
% canonical-source-hash: fnv1a64:ad2411302494231f
% canonical-lessons: IT-C352-impatto, IT-C352-ripercussione, IT-C352-risvolto, IT-C352-scatenare, IT-C352-innescare

\chapter{Le ripercussioni}
\label{ch:it-le-ripercussioni}
\hlchaptermodality{352}

\begin{chapteropening}
\textbf{By the end of this chapter:} \emph{I can describe the impact of an event, its knock-on effects and hidden sides, and say what set it off.}

Complete the last lesson of chapter 352: I can describe the impact of an event, its knock-on effects and hidden sides, and say what set it off.
\end{chapteropening}

\section[l'impatto]{l'impatto — the impact}

\label{lesson:IT-C352-impatto}



\begin{quote}
Before the new one: say the Italian for to determine, then the Italian for to influence.
\end{quote}

\subsection*{You'll want to know: l'impatto}
\textbf{l'impatto} — \textquotedblleft{}the impact\textquotedblright{}.

\begin{quote}
Il turismo ha un forte impatto sulla città.
\end{quote}

\begin{quote}
\emph{Tourism has a strong impact on the city.}
\end{quote}

Also a physical crash: \emph{l'impatto con il muro}.

\begin{cousinweb}
\emph{impatto} was borrowed from English \emph{impact} (EtymDB), which goes back to Latin \emph{impāctus}, past participle of \emph{impingere} \textquotedblleft{}to push or strike against\textquotedblright{} (Lewis \& Short), built on \emph{pangere} \textquotedblleft{}to fix, to drive in\textquotedblright{}.
\end{cousinweb}

\begin{grammarlens}[title={What you've built}]
One more word for telling, explaining and describing at B1.
\end{grammarlens}

\subsection*{Your turn}
\begin{itemize}
\raggedright
  \item \emph{Say it:} \emph{l'impatto}
  \item \emph{Say it:} \emph{l'impatto}, once more
  \item \emph{Say it:} a sentence of your own with \emph{l'impatto}
  \item \emph{Recall:} say the Italian for to determine, then the Italian for to influence, then say \emph{l'impatto} again
\end{itemize}

\begin{tcolorbox}[breakable,colback=teal!4,colframe=teal!35!black,title={Before you move on}]
What does l'impatto mean? (\textquotedblleft{}The impact\textquotedblright{}.) Say it once more.
\end{tcolorbox}
\section[la ripercussione]{la ripercussione — the repercussion, the knock-on effect}

\label{lesson:IT-C352-ripercussione}



\begin{quote}
Before the new one: say the Italian for to influence, then the Italian for the impact.
\end{quote}

\subsection*{You'll want to know: la ripercussione}
\textbf{la ripercussione} — \textquotedblleft{}the repercussion, the knock-on effect\textquotedblright{}.

\begin{quote}
La crisi ha avuto ripercussioni su tutto il paese.
\end{quote}

\begin{quote}
\emph{The crisis had repercussions on the whole country.}
\end{quote}

Mostly plural: \emph{avere ripercussioni su}.

\begin{cousinweb}
\emph{ripercussione} is a learned borrowing from Latin \emph{repercussiō} \textquotedblleft{}a striking back, a rebound\textquotedblright{} (Lewis \& Short), from \emph{repercutere} \textquotedblleft{}to strike back\textquotedblright{}, built on \emph{percutere} and \emph{quatere} \textquotedblleft{}to shake, to strike\textquotedblright{}.
\end{cousinweb}

\begin{grammarlens}[title={What you've built}]
One more word for telling, explaining and describing at B1.
\end{grammarlens}

\subsection*{Your turn}
\begin{itemize}
\raggedright
  \item \emph{Say it:} \emph{la ripercussione}
  \item \emph{Say it:} \emph{la ripercussione}, once more
  \item \emph{Say it:} a sentence of your own with \emph{la ripercussione}
  \item \emph{Recall:} say the Italian for to influence, then the Italian for the impact, then say \emph{la ripercussione} again
\end{itemize}

\begin{tcolorbox}[breakable,colback=teal!4,colframe=teal!35!black,title={Before you move on}]
What does la ripercussione mean? (\textquotedblleft{}The repercussion, the knock-on effect\textquotedblright{}.) Say it once more.
\end{tcolorbox}
\section[il risvolto]{il risvolto — the implication, the side effect; the turn-up (of trousers)}

\label{lesson:IT-C352-risvolto}



\begin{quote}
Before the new one: say the Italian for the impact, then the Italian for the repercussion.
\end{quote}

\subsection*{You'll want to know: il risvolto}
\textbf{il risvolto} — \textquotedblleft{}the implication, the side effect; the turn-up (of trousers)\textquotedblright{}.

\begin{quote}
La storia ha un risvolto che nessuno immaginava.
\end{quote}

\begin{quote}
\emph{The story has an implication that no one imagined.}
\end{quote}

The first sense is the folded-back part of a sleeve or of trousers; then the hidden side of a matter.

\begin{cousinweb}
\emph{risvolto} is built from \emph{ri-} \textquotedblleft{}back\textquotedblright{} and \emph{svolto}, a form of \emph{svolgere} and \emph{volgere} \textquotedblleft{}to turn\textquotedblright{}, which is inherited from Latin \emph{volvere} \textquotedblleft{}to roll, to turn\textquotedblright{} (Lewis \& Short): the part turned back.
\end{cousinweb}

\begin{grammarlens}[title={What you've built}]
One more word for telling, explaining and describing at B1.
\end{grammarlens}

\subsection*{Your turn}
\begin{itemize}
\raggedright
  \item \emph{Say it:} \emph{il risvolto}
  \item \emph{Say it:} \emph{il risvolto}, once more
  \item \emph{Say it:} a sentence of your own with \emph{il risvolto}
  \item \emph{Recall:} say the Italian for the impact, then the Italian for the repercussion, then say \emph{il risvolto} again
\end{itemize}

\begin{tcolorbox}[breakable,colback=teal!4,colframe=teal!35!black,title={Before you move on}]
What does il risvolto mean? (\textquotedblleft{}The implication, the side effect; the turn-up (of trousers)\textquotedblright{}.) Say it once more.
\end{tcolorbox}
\section[scatenare]{scatenare — to unleash, to spark off}

\label{lesson:IT-C352-scatenare}



\begin{quote}
Before the new one: say the Italian for the repercussion, then the Italian for the implication.
\end{quote}

\subsection*{You'll want to know: scatenare}
\textbf{scatenare} — \textquotedblleft{}to unleash, to spark off\textquotedblright{}.

\begin{quote}
La sua frase ha scatenato una lunga discussione.
\end{quote}

\begin{quote}
\emph{His remark sparked off a long discussion.}
\end{quote}

Literally to unchain. \emph{Scatenarsi} — to break loose; on a dance floor, to let yourself go.

\begin{cousinweb}
\emph{scatenare} is built from \emph{s-} \textquotedblleft{}off\textquotedblright{} and \emph{catena} \textquotedblleft{}chain\textquotedblright{}, which is inherited from Latin \emph{catēna} \textquotedblleft{}a chain\textquotedblright{} (Lewis \& Short): to let something off the chain.
\end{cousinweb}

\begin{grammarlens}[title={What you've built}]
One more word for telling, explaining and describing at B1.
\end{grammarlens}

\subsection*{Your turn}
\begin{itemize}
\raggedright
  \item \emph{Say it:} \emph{scatenare}
  \item \emph{Say it:} \emph{scatenare}, once more
  \item \emph{Say it:} a sentence of your own with \emph{scatenare}
  \item \emph{Recall:} say the Italian for the repercussion, then the Italian for the implication, then say \emph{scatenare} again
\end{itemize}

\begin{tcolorbox}[breakable,colback=teal!4,colframe=teal!35!black,title={Before you move on}]
What does scatenare mean? (\textquotedblleft{}To unleash, to spark off\textquotedblright{}.) Say it once more.
\end{tcolorbox}
\section[innescare]{innescare — to trigger, to set off}

\label{lesson:IT-C352-innescare}



\begin{quote}
Before the new one: say the Italian for the implication, then the Italian for to unleash.
\end{quote}

\subsection*{You'll want to know: innescare}
\textbf{innescare} — \textquotedblleft{}to trigger, to set off\textquotedblright{}.

\begin{quote}
Un piccolo errore può innescare un grande problema.
\end{quote}

\begin{quote}
\emph{A small mistake can trigger a big problem.}
\end{quote}

Literally to bait a hook or to prime a charge.

\begin{cousinweb}
\emph{innescare} is built from \emph{in-} and \emph{esca} \textquotedblleft{}bait; tinder\textquotedblright{}, which is inherited from Latin \emph{ēsca} \textquotedblleft{}food, bait\textquotedblright{} (Lewis \& Short): to put the bait or the spark in place.
\end{cousinweb}

\begin{grammarlens}[title={What you've built}]
One more word for telling, explaining and describing at B1.
\end{grammarlens}

\subsection*{Your turn}
\begin{itemize}
\raggedright
  \item \emph{Say it:} \emph{innescare}
  \item \emph{Say it:} \emph{innescare}, once more
  \item \emph{Say it:} a sentence of your own with \emph{innescare}
  \item \emph{Recall:} say the Italian for the implication, then the Italian for to unleash, then say \emph{innescare} again
\end{itemize}

\begin{tcolorbox}[breakable,colback=teal!4,colframe=teal!35!black,title={Before you move on}]
What does innescare mean? (\textquotedblleft{}To trigger, to set off\textquotedblright{}.) Say it once more.
\end{tcolorbox}
